Work in the frame $(k, m_1, \dots, m_{n-1})$ of Remark \ref{rem:coords} along a transversal with parameter $\tau$, with Gram $G(\tau) = \begin{psmallmatrix} \lambda & b^{\mathsf T} \\ b & G_1 \end{psmallmatrix}$, $\lambda(0) = 0$, $b(0) = 0$, $G_1(0)$ invertible, and $\lambda'(0) =: T \neq 0$, so that $\delta(k \otimes k) = T\, d\tau$. The inverse cometric $G^{-1}$, a symmetric form on $TM$, has matrix $C := [\text{Gram}]^{-1}$ with respect to the covector basis, meaning $G^{-1} = \sum C_{ab}\, \kappa_a \otimes \kappa_b$ over the frame covectors $\kappa = (k, m)$. Block inversion:
\[
C_{00} = (\lambda - b^{\mathsf T} G_1^{-1} b)^{-1} = \frac{1}{T\tau}\bigl(1 + O(\tau)\bigr), \qquad
C_{0i} = -C_{00}\, (b^{\mathsf T} G_1^{-1})_i = O(1),
\]
since $b = O(\tau)$ makes $b^{\mathsf T} G_1^{-1} b = O(\tau^2)$ and $C_{00}\,O(\tau) = O(1)$; and the $m$-block of $C$ is $G_1^{-1} + G_1^{-1} b\, C_{00}\, b^{\mathsf T} G_1^{-1} = O(1)$. Multiplying by $\tau$ and letting $\tau \to 0$, every entry goes to zero except the top-left:
\[
\tau\, G^{-1} \;\longrightarrow\; \frac{1}{T}\, k \otimes k .
\]
Independence of the extension of $k$: by Lemma \claimref{sheaf.invariant-form-t3}(2) the derivative $T = dg_k(\partial_\tau)|_\Sig$ is extension-independent. Dependence on the scale of $k$: $k \mapsto ck$ multiplies the numerator $k \otimes k$ by $c^2$ and $T$ by $c^2$ (Lemma \claimref{sheaf.invariant-form-t3}(3)); the limit is unchanged, so the residue does not depend on $k$ at all, only on the line $\shL$ and the degeneration. Dependence on $\tau$: replacing $\tau$ by $a\tau$ ($a > 0$ smooth) multiplies the limit by $a|_\Sig$, exactly the transformation of the $d\tau$-coefficient of a section of $N\Sig$; the assignment $d\tau \mapsto \lim \tau G^{-1}$ is therefore a well-defined homomorphism $N^*\Sig \to \shL^{\otimes 2}$, and the computation identifies it as $d\tau \mapsto (1/T)\,k \otimes k = \delta^{-1}(d\tau)$: the analytic residue \emph{is} the algebraic transversality datum, inverted. It is nowhere vanishing because $\delta$ is an isomorphism under (T3). A nowhere-vanishing homomorphism between line bundles is an isomorphism, giving $N^*\Sig \cong \shL^{\otimes 2}$ (equivalently, the nondegenerate pairing $(\shL^* \otimes \shL^*) \otimes N^*\Sig \to \R$ induced by $\delta^{-1}$, whence $N^*\Sig \cong (\shL^* \otimes \shL^*)^* = \shL^{\otimes 2}$); dualising gives $N\Sig \cong (\shL^*)^{\otimes 2}$. The partner identification: in the ambient completion, $k^* = \nO/(2\nu)$ satisfies $kk^* + k^*k = 1$ (Corollary \claimref{sheaf.localised-transport}), so rescaling $k \mapsto ck$ forces $k^* \mapsto k^*/c$: the partner line is canonically $\shL^*$, hence a square root of $N\Sig$. Finally, $\shL^{\otimes 2} \cong N^*\Sig$ says exactly that $k$ carries the weight of a half-power of a defining function, and $k/\sqrt t$ is weightless: the rescaling of Proposition \claimref{sheaf.exceptional-algebra-deformation-normal-cone} is the canonical normalisation.
